\section{Experiment Results}

This section reports the first retrieval results from Gemini Embedding 2. All results use cosine similarity over the generated 768-dimensional vectors. For each query, the system ranks all 533 catalogue products and evaluates whether any labelled relevant product appears in the top ranked results.

\subsection{Primary Retrieval Results}

The main table includes both matched retrieval settings and cross-representation ablations. In the matched settings, each query is compared with the product representation of the same modality. In the ablations, single-modality text or image queries are compared against product multimodal embeddings.

\begin{center}
\scriptsize
\begin{tabular}{p{0.42\textwidth}rrrrr}
\hline
\textbf{Setting} & \textbf{Queries} & \textbf{Hit@1} & \textbf{Hit@5} & \textbf{Hit@10} & \textbf{MRR} \\
\hline
Text query $\rightarrow$ product text & 20 & 0.6500 & 0.9000 & 0.9000 & 0.7308 \\
Text query $\rightarrow$ product multimodal & 20 & 0.6000 & 0.8000 & 0.9000 & 0.6913 \\
\hline
Cropped image query $\rightarrow$ product image & 11 & 0.7273 & 0.8182 & 0.8182 & 0.7576 \\
Whole image query $\rightarrow$ product image & 11 & 0.4545 & 0.6364 & 0.7273 & 0.5584 \\
\hline
Cropped image query $\rightarrow$ product multimodal & 11 & 0.7273 & 0.7273 & 0.7273 & 0.7273 \\
Whole image query $\rightarrow$ product multimodal & 11 & 0.6364 & 0.7273 & 0.7273 & 0.6667 \\
\hline
Whole image + text query $\rightarrow$ product multimodal & 11 & 0.4545 & 0.6364 & 0.8182 & 0.5707 \\
\hline
\end{tabular}
\end{center}

The text-to-text result is the strongest matched text setting. Direct metadata-like language, such as brand and ABV queries, is handled well, while flavour queries are more difficult because they require broader semantic matching across tasting notes. The text-to-multimodal ablation is slightly weaker than text-to-text, suggesting that adding product-side image evidence does not help purely textual queries in this setup. For image queries, cropped query images are easier than whole-scene query images. Product multimodal embeddings preserve cropped-image Hit@1 and improve whole-image Hit@1 compared with product image embeddings, but the whole image + text setting does not outperform whole-image retrieval in this first run.

\subsection{Text Query Breakdown}

\begin{center}
\begin{tabular}{lrrrr}
\hline
\textbf{Query style} & \textbf{Hit@1} & \textbf{Hit@5} & \textbf{Hit@10} & \textbf{MRR} \\
\hline
Direct brand lookup & 1.0000 & 1.0000 & 1.0000 & 1.0000 \\
ABV lookup & 0.7500 & 1.0000 & 1.0000 & 0.8750 \\
Semantic brand description & 0.6000 & 0.8000 & 0.8000 & 0.6400 \\
Semantic flavour & 0.3333 & 0.8333 & 0.8333 & 0.4861 \\
\hline
\end{tabular}
\end{center}

The direct brand queries are solved perfectly because the query wording aligns closely with catalogue metadata. The larger gap appears in flavour-oriented queries: relevant products often share broad sensory concepts, but not always the same exact wording or product style. This makes the flavour set a better stress test of semantic retrieval than brand lookup.

\subsection{Image Query Breakdown}

\begin{center}
\begin{tabular}{lrrrr}
\hline
\textbf{Query style} & \textbf{Hit@1} & \textbf{Hit@5} & \textbf{Hit@10} & \textbf{MRR} \\
\hline
Cropped bottle image & 0.7273 & 0.8182 & 0.8182 & 0.7576 \\
Whole scene image & 0.4545 & 0.6364 & 0.7273 & 0.5584 \\
\hline
\end{tabular}
\end{center}

The difference between cropped-bottle and whole-scene queries confirms that the visual benchmark is sensitive to query cleanliness. A product-centric crop often preserves label shape, bottle silhouette, and packaging cues. A whole scene contains the same target whisky but adds visual distractors, so the embedding must infer which object is relevant before matching it to catalogue products.

\subsection{Current Takeaways}

These results support three early observations. First, embedding-based retrieval is appropriate for this catalogue because the model can retrieve relevant products without explicit filtering or hand-written ranking rules. Second, query modality matters: visual retrieval is much stronger when the query image isolates the bottle. Third, multimodal product embeddings are not uniformly better; they help image queries at the very top of the ranking, but they do not improve text queries in this first evaluation.

\subsection{Summary}

Overall, the experiment shows that the whisky catalogue can be treated as a product-level similarity search problem. The catalogue provides enough structured text and bottle imagery for a vector embedding model to retrieve relevant products across text, image, and image-text queries. The strongest and most stable case is text retrieval, especially when the query aligns with explicit metadata such as brand or ABV. Semantic flavour queries are harder, but still often retrieve relevant products within the top five or top ten results.

The image results show a clear distinction between clean product-centric queries and realistic scene-level queries. Cropped bottle images are substantially easier because the query focuses on label and bottle appearance. Whole-scene images are more difficult because the target whisky must be identified among irrelevant visual context. This confirms that the image benchmark is not only testing product matching, but also the robustness of the embedding model to noisy visual inputs.

The multimodal results are mixed rather than uniformly better. Product multimodal embeddings improve whole-image retrieval at the top rank, which suggests that pairing product text with product imagery can make visual search more robust. However, text queries perform better against product text embeddings than against product multimodal embeddings, and whole image + text queries do not yet outperform whole image queries. The final story is therefore not that multimodal embeddings always dominate, but that their benefit depends on the query modality and the amount of noise in the query input.

In practical terms, the current evidence supports a simple retrieval design: use text embeddings for ordinary text search, use image or multimodal product embeddings for visual search, and keep cropped-image and whole-scene-image evaluation separate. This gives a clearer measurement of when multimodal embeddings add value and avoids overstating their advantage.
